\documentclass[11pt,a4paper]{article}
\usepackage{fontspec,amsmath,amssymb,amsthm,microtype}
\setmainfont{lmroman10-regular.otf}[BoldFont=lmroman10-bold.otf,ItalicFont=lmroman10-italic.otf,BoldItalicFont=lmroman10-bolditalic.otf]
\usepackage[margin=25mm]{geometry}
\usepackage[unicode,hidelinks]{hyperref}
\hypersetup{pdfauthor={LangevinLean Project}}
\setlength{\emergencystretch}{3em}
\setlength{\parskip}{0.4em}
\setlength{\parindent}{0pt}
\allowdisplaybreaks[2]
\raggedbottom
\title{\texorpdfstring{Entropy convergence from regular density interpolations}{Entropy convergence from regular density interpolations}}
\author{LangevinLean Project}
\date{}
\begin{document}\maketitle
Draft mathematical exposition; these notes may contain mistakes.


Suppose a positive regular density path represents each Gaussian ULA interpolation. Assume its weak equation admits the entropy test, its score and error energies are finite, and the stated numerical error and gradient-moment estimates hold. These assumptions imply the entropy derivative estimate, the one-step recurrence, and the iterated KL bound. The construction of the density path and the analytic estimates remain hypotheses of this conditional theorem.


\section{\texorpdfstring{Derivation from the regular interpolation assumptions}{Derivation from the regular interpolation assumptions}}


The equality between the density measure and the ULA interpolation proves normalization. Local domination of the differentiated entropy integrand gives finite KL, continuity, and the right derivative

\[H'_k(t)=\int\dot r_{k,t}(\log r_{k,t}+1)\,d\nu.\]

The finite score energy gives finite Fisher information. The admissible weak equation and Young's inequality imply \(H'_k(t)\le-3I(t)/4+E(t)\). Insert the assumed bound \(E(t)\le t^2L^2G_k+2tdL^2\) and the entropy-dependent bound on \(G_k\). The one-step comparison gives the recurrence with coefficient six, and geometric iteration gives the final coefficient eight. This derives the entropy derivative and recurrence from the listed assumptions; it does not derive those assumptions merely from LSI.



Earlier results used here: \cite[Section~3.26.3]{LangevinLean2026ULASC}.

\begin{thebibliography}{ULA-LSI}
\bibitem[ULA-SC]{LangevinLean2026ULASC}
LangevinLean Project. \emph{Unadjusted Langevin under Strong Convexity}. LL-01, 2026. Draft report, content revision 7e430ee07530. \url{https://odakyildiz.github.io/LangevinLean/docs/pdf/UnadjustedLangevin/FullProof.pdf}.
\end{thebibliography}

\end{document}
